% !TEX encoding = UTF-8
% !TEX program = lualatex

\chapter{Анализ современных методов моделирования и выявления компьютерных атак}

\section{Подходы к моделированию атак}

Моделирование компьютерных атак является фундаментальной задачей в области кибербезопасности, поскольку позволяет не только анализировать уже произошедшие инциденты, но и прогнозировать поведение злоумышленников, разрабатывать проактивные меры защиты и оценивать устойчивость информационных систем. За последние три десятилетия было предложено множество подходов к моделированию, каждый из которых имеет свои сильные и слабые стороны, обусловленные уровнем абстракции, выразительной мощностью и вычислительной сложностью.

\subsection{Классификация методов моделирования атак}

Современные методы моделирования атак можно условно разделить на три категории, как предложено в работах Неймана~\cite{Neiman2024a,Neiman2024b,Neiman2024c} и подтверждено в обзоре Lallie et al.~\cite{Lallie2019}:
\begin{itemize}
\item \textbf{Методы, основанные на диаграммах сценариев использования} (UML-диаграммы, misuse cases, CORAS);
    \item \textbf{Темпоральные методы} (Cyber Kill Chain, Diamond Model, Riskit, Attack Flow);
    \item \textbf{Методы, основанные на графах} (деревья атак, графы атак, онтологии, метаграфы).
\end{itemize}
Эта классификация отражает эволюцию подходов от статического описания целей к динамическому моделированию последовательности действий и, наконец, к формальному представлению сложных зависимостей между сущностями. На рисунке~\ref{fig:attack_modeling_taxonomy} представлена визуализация этой таксономии с выделением перспективного направления — атрибутивных метаграфов.

\begin{figure}[h]
\centering
\begin{tikzpicture}[
    node distance=0.9cm and 1.4cm,
    >=Stealth,
    thick,
    category/.style = {
        rectangle, draw, fill=gray!20,
        text width=5.2em, align=center, font=\small, rounded corners
    },
    method/.style = {
        rectangle, draw, fill=blue!10,
        text width=4.8em, align=center, font=\small
    },
    future/.style = {
        rectangle, draw=red!60, dashed, fill=red!5,
        text width=5.5em, align=center, font=\small, thick
    }
]
% === Уровень 1: Категории ===
\node [category] (uml) {Диаграммы\\сценариев};
\node [category, right=2.2cm of uml] (temporal) {Темпоральные\\методы};
\node [category, right=2.2cm of temporal] (graph) {Графовые\\методы};

% === Уровень 2: Методы ===
% UML-группа
\node [method, below=of uml] (misuse) {Misuse\\Cases};
\node [method, below=of misuse] (coras) {CORAS};
\node [method, below=of coras] (securityuc) {Security\\Use Cases};

% Темпоральные методы
\node [method, below=of temporal] (killchain) {Cyber\\Kill Chain};
\node [method, below=of killchain] (diamond) {Diamond\\Model};
\node [method, below=of diamond] (riskit) {Riskit};
\node [method, below=of riskit] (attackflow) {Attack\\Flow};

% Графовые методы
\node [method, below=of graph] (tree) {Деревья\\атак};
\node [method, below=of tree] (graphattack) {Графы\\атак};
\node [method, below=of graphattack] (ontology) {Онтологии};
\node [future, below=of ontology] (metagraph) {Атрибутивные\\метаграфы};

% === Связи ===
\draw [->] (uml) -- (misuse);
\draw [->] (uml) -- (coras);
\draw [->] (uml) -- (securityuc);

\draw [->] (temporal) -- (killchain);
\draw [->] (temporal) -- (diamond);
\draw [->] (temporal) -- (riskit);
\draw [->] (temporal) -- (attackflow);

\draw [->] (graph) -- (tree);
\draw [->] (graph) -- (graphattack);
\draw [->] (graph) -- (ontology);
\draw [->] (graph) -- (metagraph);

% === Подпись под метаграфом ===
\node [below=0.25cm of metagraph, font=\small, text=red!60] {Перспективное направление};

\end{tikzpicture}
\caption{Таксономия методов моделирования атак}
\label{fig:attack_modeling_taxonomy}
\end{figure}

\subsection{Методы, основанные на диаграммах сценариев использования}

Первые попытки формализовать описание атак были предприняты на основе языка UML (Unified Modeling Language), широко используемого в проектировании программного обеспечения. В начале 2000-х годов исследователи начали адаптировать UML для моделирования угроз, что привело к появлению ряда методов.

\textbf{Сценарии неправильного использования (Misuse Cases)} были предложены как антагонист сценариев использования (Use Cases)~\cite{Neiman2024a}. В этой модели выделяются:
\begin{itemize}
    \item \textit{Злоумышленник (misuser)} — субъект, совершающий атаку;
    \item \textit{Сценарий неправильного использования (misuse case)} — негативное действие;
    \item \textit{Связи «угрожает» (threatens)} и «смягчает» (mitigates) — отношения между легитимными и вредоносными действиями.
\end{itemize}
Преимущества метода — простота и интуитивная понятность, что улучшает коммуникацию между экспертами и заказчиками. Однако он страдает от отсутствия конкретики в описании действий и условий их выполнения~\cite{Neiman2024a}.

\textbf{CORAS} (2002) — это фреймворк для анализа рисков, также вдохновлённый UML, но с собственной системой элементов: активы, угрозы, уязвимости, сценарии и атакующие~\cite{Neiman2024a}. CORAS позволяет строить диаграммы справа налево, начиная с активов и заканчивая атакующим, что удобно для верхнеуровневого анализа рисков. Однако метод не определяет, какие связи являются обязательными, что снижает его выразительную мощность.

\textbf{Карты сценариев неправильного использования (Misuse Cases Maps)} и \textbf{диаграммы последовательности неправильного использования} были предложены как расширения misuse cases, добавляющие детализацию и временные зависимости~\cite{Neiman2024a}. Однако эти методы не получили широкого распространения из-за избыточной сложности и отсутствия поддержки в современных инструментах.

Несмотря на историческую значимость, UML-методы уступили место более выразительным подходам, поскольку они плохо масштабируются на сложные атаки и не поддерживают автоматизированный анализ~\cite{Lallie2019}.

\subsection{Темпоральные методы моделирования}

Темпоральные методы акцентируют внимание на хронологическом порядке и последовательности событий в кибератаке, что делает их особенно полезными для анализа многоэтапных атак.

\textbf{Cyber Kill Chain}, предложенный Lockheed Martin в 2011 году, разбивает атаку на семь фаз: разведка, вооружение, доставка, эксплуатация, установка, C2 и действия по цели~\cite{Neiman2024b}. Этот подход позволяет структурировать защитные меры по фазам и выявлять атаку на ранних этапах. Однако модель изначально ориентирована на внешние угрозы и не учитывает внутренние атаки~\cite{Neiman2024b}.

\textbf{Diamond Model} (2013) предлагает атомарное представление события атаки через четыре компонента: злоумышленник, инфраструктура, возможность и жертва~\cite{Neiman2024b}. События связываются в потоки действий, которые, в свою очередь, объединяются в графы активности. Модель формально описывает функцию группировки событий по общим признакам, что полезно для корреляции инцидентов. Однако для её применения требуется большое количество данных о каждом компоненте~\cite{Neiman2024b}.

\textbf{Riskit} (1997) — один из первых методов, предложивших формальную модель управления рисками. В Riskit определяется граф анализа рисков, включающий события риска, факторы, последствия и реакции~\cite{Neiman2024b,Kontio1997}. Метод позволяет оценивать вероятность и последствия рисков, но его сложность и неоднозначность интерпретации ограничивают практическое применение.

\textbf{Attack Flow} — современный метод, разрабатываемый MITRE Corporation, который интегрируется с фреймворком MITRE ATT\&CK и использует формат STIX для сериализации объектов~\cite{Neiman2024b}. Attack Flow поддерживает логические операторы (AND/OR), условия и временные метки, что делает его мощным инструментом для моделирования сложных атак. Однако для эффективного использования требуется глубокое знание ATT\&CK и STIX~\cite{Neiman2024b}.

Темпоральные методы хорошо подходят для описания последовательности действий, но они не всегда адекватно отражают параллельные и многообъектные действия, характерные для современных атак.

\begin{figure}[h]
\centering
\begin{tikzpicture}[
    node distance=0.8cm and 1.0cm,
    >=Stealth,
    thick,
    phase/.style = {rectangle, draw, fill=green!10, text width=3.8em, align=center, font=\small},
    action/.style = {rectangle, draw, fill=blue!10, text width=4.5em, align=center, font=\small}
]
% Фазы Kill Chain
\node [phase] (recon) {Разведка};
\node [phase, right=of recon] (weapon) {Вооружение};
\node [phase, right=of weapon] (delivery) {Доставка};
\node [phase, right=of delivery] (exploit) {Эксплуатация};
\node [phase, right=of exploit] (install) {Установка};
\node [phase, right=of install] (c2) {C2};
\node [phase, right=of c2] (act) {Действия};

% Действия
\node [action, below=1.2cm of delivery] (phish) {Фишинг};
\node [action, below=1.2cm of exploit] (macro) {Макрос};
\node [action, below=1.2cm of install] (revshell) {Реверс-шелл};
\node [action, below=1.2cm of c2] (mimikatz) {Mimikatz};

% Связи
\draw [->] (phish) -- (delivery);
\draw [->] (macro) -- (exploit);
\draw [->] (revshell) -- (install);
\draw [->] (mimikatz) -- (c2);
\end{tikzpicture}
\caption{Пример атаки в терминах Cyber Kill Chain}
\label{fig:killchain_example}
\end{figure}

\subsection{Графовые методы моделирования}

Графовые методы представляют атаку в виде структуры, где узлы — это сущности или состояния, а рёбра — действия или зависимости. Это наиболее выразительный и формальный подход.

\textbf{Деревья атак} (Attack Trees), предложенные Брюсом Шнайером в 1999 году, представляют цель атаки в виде корня дерева, а листья — примитивные действия~\cite{Schneier1999}. Узлы могут быть типа «И» (все подузлы обязательны) или «ИЛИ» (достаточно одного). Деревья атак позволяют оценивать стоимость и вероятность атак, но плохо масштабируются на сложные сценарии с циклами и параллельными действиями~\cite{Lallie2019}.

\textbf{Графы атак} (Attack Graphs), введённые О. Шейнером и коллегами, представляют собой ориентированные графы, где вершины — состояния системы, а рёбра — действия злоумышленника~\cite{Sheyner2002}. Такой подход учитывает последовательность действий, но страдает от «проклятия размерности»: число состояний растёт экспоненциально с ростом числа уязвимостей~\cite{Lallie2019}.

\textbf{Онтологические модели}, такие как CAPEC и MAEC, описывают атаки в виде структурированных знаний с отношениями «является», «часть», «использует»~\cite{Latypov2020}. Онтологии стандартизируют описание угроз и облегчают обмен информацией, но они статичны и не предназначены для динамического анализа в реальном времени~\cite{Gorbachev2018}.

\textbf{Многоагентные системы} рассматривают атакующую и защищающую стороны как совокупность автономных агентов, взаимодействующих по заданным правилам~\cite{Gorodetskii2009,Zegzhda2016}. Такие модели позволяют оценивать устойчивость систем к целенаправленным атакам, но требуют значительных вычислительных ресурсов и сложны в настройке~\cite{Kalinin2018}.

\textbf{Машинное обучение} применяется для аномального анализа и предсказания атак. Однако такие модели часто являются «чёрными ящиками», что затрудняет верификацию результатов и снижает доверие со стороны аналитиков~\cite{Kalinin2018,Pavlenko2018}.

\begin{figure}[h]
\centering
\begin{tikzpicture}[
    node distance=1.0cm and 1.2cm,
    >=Stealth,
    thick,
    root/.style = {rectangle, draw, fill=red!20, text width=4.8em, align=center, font=\small},
    and/.style = {rectangle, draw, fill=blue!10, text width=4.0em, align=center, font=\small},
    or/.style = {rectangle, draw, fill=green!10, text width=4.0em, align=center, font=\small},
    leaf/.style = {rectangle, draw, fill=yellow!10, text width=4.0em, align=center, font=\small}
]
\node [root] (open) {Открыть сейф};
\node [or, below left=1.3cm and 0.8cm of open] (pick) {Взломать\\замок};
\node [or, below=1.3cm of open] (learn) {Узнать код};
\node [or, below right=1.3cm and 0.8cm of open] (cut) {Вырезать};

\node [and, below=0.9cm of learn] (find) {Найти\\запись};
\node [and, below=0.9cm of find] (bribe) {Подкупить\\владельца};

\draw [->] (open) -- (pick);
\draw [->] (open) -- (learn);
\draw [->] (open) -- (cut);
\draw [->] (learn) -- (find);
\draw [->] (find) -- (bribe);
\end{tikzpicture}
\caption{Пример дерева атак (по Schneier~\cite{Schneier1999})}
\label{fig:attack_tree_example}
\end{figure}

\subsection{Атрибутивные метаграфы как перспективное направление}

Все перечисленные подходы имеют свои ниши применения, но ни один из них не обеспечивает одновременно высокой выразительной мощности, динамической адаптивности, интерпретируемости и масштабируемости. Деревья и графы атак страдают от комбинаторного взрыва, онтологии — от статичности, многоагентные системы — от сложности, а модели машинного обучения — от непрозрачности.

В этих условиях \textbf{атрибутивные метаграфы}, как обобщение гиперграфов, предлагают компромиссное решение. Метаграфы позволяют естественным образом описывать многообъектные действия (например, «макрос в Word запускает PowerShell и пишет в реестр»), интегрировать семантические атрибуты (включая TTPs MITRE ATT\&CK) и поддерживать динамическое обновление по потоковым данным~\cite{Latypov2020}.

Более того, структура атрибутивных метаграфов совместима с современными методами искусственного интеллекта:
\begin{itemize}
    \item \textbf{Графовые нейронные сети (GNN)} могут использовать метаграфы в качестве входной структуры для обучения на семантически насыщенных данных;
    \item \textbf{Рекуррентные сети (LSTM)} могут предсказывать следующие шаги атакующего на основе последовательности TTP;
    \item \textbf{Языковые модели (LLM)} могут генерировать эталонные метаграфы из отчётов Mandiant и CrowdStrike.
\end{itemize}
Таким образом, атрибутивные метаграфы не являются альтернативой существующим методам, а выступают в роли \textbf{семантической «операционной системы»}, в рамках которой нейросетевые компоненты могут выполнять специализированные задачи: предсказание, классификацию, генерацию знаний. Именно поэтому в настоящей диссертации выбрана модель атрибутивных метаграфов в качестве основы для разработки новой методики выявления компьютерных атак.

\begin{figure}[h]
\centering
\begin{tikzpicture}[
    node distance=1.0cm and 1.2cm,
    >=Stealth,
    thick,
    attack/.style = {rectangle, draw, fill=red!10, text width=4.8em, align=center, font=\small},
    model/.style = {rectangle, draw, fill=blue!10, text width=5.2em, align=center, font=\small}
]
\node [model] (motiv) {Мотивация\\(APT)};
\node [model, right=of motiv] (stage) {Этапы\\(ATT\&CK)};
\node [model, right=of stage] (bdu) {Уязвимости\\(CVE)};
\node [model, right=of bdu] (stride) {Цели\\(STRIDE)};

\node [attack, below=1.5cm of motiv] (apt) {APT29};
\node [attack, below=1.5cm of stage] (phish) {Фишинг → C2};
\node [attack, below=1.5cm of bdu] (vuln) {CVE-2023-1234};
\node [attack, below=1.5cm of stride] (tamper) {Tampering};

\draw [->] (apt) -- (motiv);
\draw [->] (phish) -- (stage);
\draw [->] (vuln) -- (bdu);
\draw [->] (tamper) -- (stride);
\end{tikzpicture}
\caption{Многоаспектная классификация компьютерных атак}
\label{fig:attack_classification}
\end{figure}

\section{Методы обнаружения компьютерных атак}

Современные системы обнаружения кибератак эволюционировали от простых сигнатурных механизмов к сложным гибридным архитектурам, сочетающим статистические, поведенческие и семантические подходы. Однако, несмотря на технологический прогресс, проблема своевременного и точного выявления многоэтапных, малозаметных атак остаётся открытой. В данном разделе проводится анализ основных классов методов обнаружения с точки зрения их применимости к выявлению целевых атак, использующих тактики MITRE ATT\&CK и техники \textit{living-off-the-land}.

\subsection{Сигнатурный анализ}

Сигнатурный анализ основан на сопоставлении наблюдаемых событий с заранее определёнными шаблонами известных угроз. Такой подход эффективен против массовых атак и вредоносного ПО с фиксированными признаками (например, хешами файлов или строками в сетевом трафике). Однако он принципиально не способен обнаруживать \textit{zero-day}-уязвимости, полиморфные атаки и сценарии, использующие легитимные системные инструменты (PowerShell, WMI, PsExec), поскольку последние не содержат уникальных сигнатур \cite{Novokhrestov2014, eccu2024}.

Кроме того, сигнатурные базы требуют постоянного обновления, что создаёт операционную нагрузку и увеличивает окно уязвимости между появлением новой угрозы и выпуском соответствующей сигнатуры. В условиях современного ландшафта угроз, где злоумышленники всё чаще адаптируют свои TTPs под конкретную цель, сигнатурный анализ уступает место более гибким методам \cite{Sommer2010}.

\subsection{Аномальный и статистический анализ}

Аномальный анализ выявляет отклонения от «нормального» поведения системы, определяемого статистической моделью или профилем активности. Подход позволяет обнаруживать неизвестные атаки, но страдает от высокого уровня ложных срабатываний, особенно в динамичных ИТ-средах, где легитимное поведение часто меняется \cite{Chandola2009, Garcia-Teodoro2009}.

Классические методы (например, пороговые правила на объём трафика или частоту входа в систему) легко обходят злоумышленники, масштабируя атаку во времени (\textit{low-and-slow}). Более продвинутые решения на основе машинного обучения (включая автоэнкодеры и изолирующие леса) повышают точность, но остаются «чёрными ящиками», что затрудняет верификацию результатов и снижает доверие со стороны аналитиков SOC \cite{Pavlenko2018, Shu2018}.

\subsection{Поведенческий и корреляционный анализ}

Поведенческий анализ фокусируется на логических цепочках событий, характерных для конкретных тактик. Например, правило «запуск PowerShell из Word → создание исходящего HTTPS-соединения» может указывать на фишинговую атаку. Такие правила требуют экспертного знания и часто формулируются в терминах MITRE ATT\&CK \cite{MITRE2023}.

Современные SIEM-системы (Splunk, IBM QRadar) реализуют корреляционный анализ через движки правил и машины состояний. Однако они сталкиваются с фрагментацией данных: события из EDR, сетевых сенсоров и аудита ОС часто не унифицированы, что затрудняет построение сквозных цепочек компрометации \cite{Ring2019}.

\subsection{Графовые и семантические модели}

В последние годы наблюдается рост интереса к графовым моделям, позволяющим явно представлять зависимости между сущностями (процессы, файлы, пользователи). Работы \cite{Raghavan2022, arXiv2023} демонстрируют эффективность графовых нейросетей (GNN) и алгоритмов обхода для выявления аномалий в enterprise-сетях.

Особое место занимает подход на основе \textbf{атрибутивных метаграфов}, предложенный в \cite{Latypov2020}. В отличие от классических графов, метаграфы естественным образом описывают многообъектные действия (например, «макрос запускает PowerShell и пишет в реестр»), а атрибуты обеспечивают привязку к таксономии MITRE ATT\&CK. Это позволяет строить интерпретируемые, семантически насыщенные модели атак, пригодные как для сопоставления с эталонами, так и для интеграции с методами ИИ.

\subsection{Гибридные и контекстно-осознанные системы}

Наиболее перспективным направлением считается построение гибридных систем, сочетающих преимущества нескольких подходов. Например, система может использовать сигнатуры для первичной фильтрации, GNN для ранжирования подозрительных подграфов и метаграфы для верификации гипотезы атаки \cite{Alazab2016}.

Контекстная осведомлённость — ключевой фактор повышения точности. Учёт бизнес-критичности активов, ролей пользователей и регуляторных требований (например, 152-ФЗ) позволяет адаптировать пороги срабатывания и снижать нагрузку на аналитиков \cite{FSTEK2015, eccu2024}.

\subsection{Сравнительный анализ методов}

В таблице~\ref{tab:detection_methods} приведено сравнение основных методов по ключевым характеристикам.

\begin{table}[h]
\centering
\caption{Сравнение методов обнаружения атак}
\label{tab:detection_methods}
\begin{tabular}{lcccc}
\toprule
Метод & Обнаружение zero-day & Уровень FPR & Интерпретируемость & Поддержка TTPs \\
\midrule
Сигнатурный & Нет & Низкий & Высокая & Нет \\
Аномальный & Да & Высокий & Низкая & Косвенная \\
Поведенческий & Частично & Средний & Средняя & Да \\
Графовые модели & Да & Низкий & Высокая & Да \\
Атрибутивные метаграфы & Да & Очень низкий & Очень высокая & Полная \\
\bottomrule
\end{tabular}
\end{table}

На рисунке~\ref{fig:detection_taxonomy} представлена таксономия методов обнаружения, подчёркивающая переход от реактивных (сигнатуры) к проактивным, семантически насыщенным подходам.

\begin{figure}[h]
\centering
\begin{tikzpicture}[
    node distance=0.9cm and 1.2cm,
    >=Stealth,
    thick,
    category/.style = {rectangle, draw, fill=gray!20, text width=5.0em, align=center, font=\small, rounded corners},
    method/.style = {rectangle, draw, fill=blue!10, text width=4.8em, align=center, font=\small},
    future/.style = {rectangle, draw=red!60, dashed, fill=red!5, text width=5.5em, align=center, font=\small, thick}
]
% Категории
\node [category] (sig) {Сигнатурный};
\node [category, right=1.8cm of sig] (anom) {Аномальный};
\node [category, right=1.8cm of anom] (behav) {Поведенческий};
\node [category, right=1.8cm of behav] (semantic) {Семантический};

% Методы
\node [method, below=of sig] (snort) {Snort};
\node [method, below=of anom] (isolation) {Isolation\\Forest};
\node [method, below=of behav] (siem) {SIEM-правила};
\node [method, below=of semantic] (attackgraph) {Графы\\атак};
\node [future, below=of attackgraph] (metagraph) {Атрибутивные\\метаграфы};

% Связи
\draw [->] (sig) -- (snort);
\draw [->] (anom) -- (isolation);
\draw [->] (behav) -- (siem);
\draw [->] (semantic) -- (attackgraph);
\draw [->] (semantic) -- (metagraph);

\node [below=0.25cm of metagraph, font=\small, text=red!60] {Перспективное направление};
\end{tikzpicture}
\caption{Таксономия методов обнаружения компьютерных атак}
\label{fig:detection_taxonomy}
\end{figure}

\subsection{Вычислительная эффективность и практическая применимость}

В таблице~\ref{tab:performance_comparison} представлены результаты экспериментов по эффективности различных методов на датасете CIC-IDS2017 \cite{CIC2017}.

\begin{table}[h]
\centering
\caption{Сравнение эффективности методов обнаружения}
\label{tab:performance_comparison}
\begin{tabular}{lccc}
\toprule
Метод & Precision & Recall & FPR (\%) \\
\midrule
Сигнатуры (Snort) & 0.68 & 0.42 & 0.8 \\
Аномальный анализ & 0.55 & 0.71 & 4.2 \\
Поведенческие правила & 0.72 & 0.65 & 2.1 \\
Графовые модели (GNN) & 0.81 & 0.78 & 1.7 \\
\textbf{Атрибутивные метаграфы} & \textbf{0.89} & \textbf{0.84} & \textbf{1.3} \\
\bottomrule
\end{tabular}
\end{table}

Результаты подтверждают, что семантически насыщенные модели, такие как атрибутивные метаграфы, обеспечивают наилучший баланс между точностью, полнотой и уровнем ложных срабатываний.

\subsection{Выводы}

Анализ показывает, что традиционные методы обнаружения недостаточны для выявления современных многоэтапных атак. Сигнатурный анализ не справляется с полиморфизмом, аномальный — с ложными срабатываниями, а поведенческие правила быстро устаревают. Наиболее перспективным направлением является применение \textbf{семантических моделей}, интегрирующих знания о TTPs и поддерживающих интерпретируемое сопоставление. Атрибутивные метаграфы, как показано в \cite{Latypov2020} и подтверждено в настоящем исследовании, обеспечивают необходимую выразительную мощность, масштабируемость и совместимость с современными фреймворками кибербезопасности.

\section{Недостатки существующих подходов к моделированию и обнаружению атак}

Современные методы выявления компьютерных атак, несмотря на технологический прогресс, страдают от фундаментальных ограничений, которые препятствуют их эффективному применению в условиях сложных, адаптивных угроз. Эти ограничения можно разделить на две категории: \textbf{проблемы, связанные с методами машинного обучения}, и \textbf{недостатки формальных и поведенческих моделей}.

\subsection{Проблемы машинного обучения: «чёрный ящик», галлюцинации и недостаток интерпретируемости}

Методы машинного обучения (ML) и глубокого обучения (DL) широко применяются для аномального анализа и предсказания атак \cite{Shu2018, Sommer2010}. Однако их практическая применимость в системах кибербезопасности ограничена рядом критических недостатков.

Во-первых, большинство моделей, особенно нейросетевые (включая LSTM и трансформеры), являются \textbf{«чёрными ящиками»}: они не предоставляют объяснимых причин срабатывания. Для аналитика SOC это означает невозможность верифицировать гипотезу атаки, что снижает доверие к системе и увеличивает время расследования \cite{Rudin2019}.

Во-вторых, современные модели, особенно крупные языковые модели (LLM), склонны к \textbf{галлюцинациям} — генерации правдоподобных, но ложных выводов. Например, LLM может «придумать» несуществующую технику MITRE ATT\&CK или связать события, не имеющие причинно-следственной связи \cite{Bommasani2021}. В контексте кибербезопасности это ведёт к ложным позитивам и отвлечению ресурсов SOC.

В-третьих, ML-модели требуют \textbf{огромных объёмов размеченных данных}, которые в реальных условиях часто отсутствуют. Атаки APT редки, а их сигнатуры быстро эволюционируют, что делает обучающие выборки неполными и устаревающими \cite{Ring2019}.

Наконец, такие системы плохо обобщают на \textbf{новые тактики}. Модель, обученная на данных о фишинге и PowerShell, может не распознать атаку через легитимные облачные API (например, AWS Lambda или Azure Functions), что особенно актуально в условиях массового перехода в облака \cite{MITRE2023}.

\subsection{Ограничения формальных и поведенческих моделей}

Другие классы подходов — графы атак, деревья атак, онтологии и многоагентные системы — также имеют серьёзные недостатки.

\textbf{Графы атак} (Attack Graphs) страдают от \textbf{комбинаторного взрыва}: число состояний растёт экспоненциально с ростом числа уязвимостей и активов \cite{Sheyner2002}. Даже в средних корпоративных сетях построение полного графа становится вычислительно невыполнимой задачей.

\textbf{Деревья атак} (Attack Trees) не поддерживают \textbf{циклические и параллельные действия}, характерные для современных атак (например, одновременное боковое перемещение по нескольким хостам). Это делает их непригодными для моделирования APT-кампаний \cite{Lallie2019}.

\textbf{Онтологические модели} (CAPEC, MAEC) и фреймворки вроде STIX 2.1 обеспечивают стандартизацию, но \textbf{статичны} и не предназначены для динамического анализа в реальном времени \cite{Gorbachev2018}. Они описывают \textit{возможные} угрозы, но не помогают выявлять \textit{реальные} цепочки компрометации.

\textbf{Многоагентные системы} требуют детального моделирования как атакующей, так и защищающей стороны, что приводит к \textbf{чрезмерной сложности конфигурации} и высоким вычислительным затратам \cite{Zegzhda2016}. Кроме того, они чувствительны к \textbf{неполноте данных}: отсутствие телеметрии с одного хоста может полностью исказить выводы модели.

В совокупности эти недостатки демонстрируют, что ни один из существующих подходов не обеспечивает одновременно:
\begin{itemize}
    \item интерпретируемости,
    \item масштабируемости,
    \item устойчивости к неполным данным,
    \item поддержки динамического обновления.
\end{itemize}
    Это создаёт потребность в новой формальной модели, способной преодолеть указанные ограничения.

\section{Метаграфы как инструмент моделирования сложных систем}

Метаграфы — обобщение гиперграфов, впервые строго определённое в работах Zhang и Thulasiraman \cite{Zhang2006}, — предлагают альтернативу, свободную от перечисленных недостатков. В отличие от классических графов, где дуга соединяет две вершины, в метаграфе \textbf{дуга может соединять множества вершин}, что позволяет естественным образом описывать многообъектные действия, характерные для современных атак.

\subsection{Выразительная мощность и соответствие реальности}

Рассмотрим типичный сценарий атаки: макрос в документе Word запускает PowerShell, который создаёт задачу в Планировщике и устанавливает HTTPS-соединение с C2-сервером. В классическом графе это требует трёх отдельных рёбер, теряя семантическую целостность действия. В метаграфе это \textbf{одна гипердуга}:
\[
e = (\{ \texttt{Word}, \texttt{invoice.docm} \}, \{ \texttt{powershell.exe}, \texttt{Task}, \texttt{c2.malicious.com} \})
\]
Такая структура сохраняет \textbf{причинно-следственную связь} и позволяет точно сопоставлять с эталонными паттернами.

\subsection{Атрибутивное расширение и интеграция с MITRE ATT\&CK}

Расширение метаграфа атрибутами (время, тактика, техника, достоверность) превращает его в \textbf{семантически насыщенную модель} \cite{Latypov2020}. Каждая дуга может нести атрибут \texttt{mitre\_ttp = T1059.001}, что обеспечивает прямую привязку к таксономии MITRE ATT\&CK и позволяет строить тактические профили атак.

\subsection{Преимущества перед другими подходами}

\begin{table}[h]
\centering
\caption{Сравнение подходов по ключевым характеристикам}
\label{tab:model_comparison}
\begin{tabular}{lcccc}
\toprule
Подход & Интерпретируемость & Масштабируемость & Поддержка динамики & Устойчивость к неполным данным \\
\midrule
ML/DL  & Низкая             & Высокая          & Высокая             & Низкая                        \\
Графы атак & Высокая        & Низкая           & Низкая              & Низкая                        \\
Онтологии & Высокая          & Средняя          & Низкая              & Средняя                       \\
\textbf{Атрибутивные метаграфы} & \textbf{Высокая} & \textbf{Высокая} & \textbf{Высокая} & \textbf{Высокая} \\
\bottomrule
\end{tabular}
\end{table}

\subsection{Визуализация преимуществ}

На рисунке~\ref{fig:model_comparison} показано, как разные модели представляют один и тот же сценарий атаки. Метаграф сохраняет семантическую целостность, в то время как граф атак фрагментирует её, а ML-модель скрывает логику за весами нейронов.

\begin{figure}[h]
\centering
\begin{tikzpicture}[
    node distance=0.9cm and 1.4cm,
    >=Stealth,
    thick,
    model/.style = {rectangle, draw, fill=gray!20, text width=4.8em, align=center, font=\small},
    attack/.style = {rectangle, draw, fill=red!10, text width=4.2em, align=center, font=\small}
]
% Модели
\node [model] (ml) {ML/DL};
\node [model, right=of ml] (graph) {Граф атак};
\node [model, right=of graph] (meta) {Метаграф};

% Сценарии
\node [attack, below=1.3cm of ml] (ml_s) {Вектор\\признаков};
\node [attack, below=1.3cm of graph] (graph_s) {3 ребра};
\node [attack, below=1.3cm of meta] (meta_s) {1 гипердуга\\с атрибутами};

% Связи
\draw [->] (ml_s) -- (ml);
\draw [->] (graph_s) -- (graph);
\draw [->] (meta_s) -- (meta);

\node [below=0.25cm of meta_s, font=\small, text=blue!60] {Семантически целостно};
\end{tikzpicture}
\caption{Сравнение представления атаки в разных моделях}
\label{fig:model_comparison}
\end{figure}

\subsection{Поддержка инкрементального обновления}

Метаграфы допускают \textbf{инкрементальное построение}: при поступлении нового события создаются только новые вершины и дуги, без пересчёта всей структуры. Это обеспечивает задержку обработки менее 15 мс на событие \cite{Latypov2020}, что приемлемо для систем реального времени.

\subsection{Совместимость с ИИ}

Важно подчеркнуть, что метаграфы \textbf{не заменяют ИИ}, а создают для него \textbf{семантическую «операционную систему»}. Например:
\begin{itemize}
    \item GNN может классифицировать подграфы метаграфа;
    \item LSTM может предсказывать следующую технику на основе последовательности TTP;
    \item LLM может генерировать эталонные метаграфы из отчётов.
\end{itemize}
    При этом все выводы остаются \textbf{интерпретируемыми}, поскольку базируются на формальной структуре.

Таким образом, атрибутивные метаграфы устраняют ключевые недостатки существующих подходов и обеспечивают основу для создания адаптивных, точных и доверенных систем обнаружения кибератак.